\documentclass[12pt]{article}

%================================================
% PAQUETES
%================================================

\usepackage[utf8]{inputenc}
\usepackage[T1]{fontenc}
\usepackage[spanish,es-noshorthands]{babel}
\usepackage{lmodern}
\usepackage{amsmath,amssymb,amsfonts}
\usepackage{enumitem}
\usepackage{geometry}
\usepackage{array}
\usepackage{booktabs}
\usepackage{multicol}
\usepackage{fancyhdr}
\usepackage{tcolorbox}
\usepackage{siunitx}

\geometry{margin=2.5cm}

\sisetup{
	output-decimal-marker={,}
}

%================================================
% ENCABEZADO
%================================================

\pagestyle{fancy}
\fancyhf{}
\fancyhead[L]{Fundamentos de Matemática Básica}
\fancyhead[R]{Guía de Ejercicios}
\fancyfoot[C]{\thepage}

%================================================
% DATOS
%================================================

\title{\textbf{Guía de Ejercicios de Nivel Superior}\\
	Modelación Matemática y Ecuaciones de Primer Grado\\
	Contextualizada a Agronomía e Ingeniería Agrícola}

\author{Prof. Luis Tulio González Alcaino}
\date{Universidad Santo Tomás \\ Departamento de Ciencias Básicas \\ 2026}

%================================================
% DOCUMENTO
%================================================

\begin{document}
	
	\maketitle
	
	\vspace{-1em}
	
	\begin{tcolorbox}[colback=green!5,colframe=black]
		\textbf{Objetivo de la guía:}
		
		Aplicar ecuaciones de primer grado enteras, fraccionarias y literales
		en situaciones reales de Agronomía e Ingeniería Agrícola,
		desarrollando capacidad de modelación, análisis crítico,
		detección de patrones e identificación de inconsistencias.
	\end{tcolorbox}
	
	%================================================
	\section*{I. Modelación Matemática}
	
	\begin{enumerate}[label=\textbf{\arabic*.}]
		
		\item Un predio agrícola necesita aplicar fertilizante nitrogenado
		a razón de 180 kg por hectárea.
		
		Si el terreno posee \(x\) hectáreas y ya existen 240 kg almacenados,
		la cantidad total necesaria debe ser 1320 kg.
		
		\begin{enumerate}[label=\alph*)]
			\item Plantee la ecuación.
			\item Determine la cantidad de hectáreas del terreno.
			\item Interprete el resultado.
		\end{enumerate}
		
		%------------------------------------------------
		
		\item Una empresa agrícola vende cajas de tomates.
		
		Cada caja contiene 25 kg y se desea completar
		una carga total de 3.500 kg.
		
		Si ya se han cargado 1.250 kg:
		
		\begin{enumerate}[label=\alph*)]
			\item Modele mediante una ecuación.
			\item Determine cuántas cajas faltan.
			\item Analice si el resultado requiere redondeo.
		\end{enumerate}
		
	\end{enumerate}
	
	%================================================
	\section*{II. Ecuaciones Enteras}
	
	\begin{enumerate}[label=\textbf{\arabic*.}, start=3]
		
		\item Resolver:
		
		\[
		18x + 240 = 960
		\]
		
		Interpretar \(x\) como número de surcos de cultivo.
		
		%------------------------------------------------
		
		\item Resolver:
		
		\[
		45x - 120 = 780
		\]
		
		Interpretar \(x\) como cantidad de sacos de semilla.
		
		%------------------------------------------------
		
		\item Resolver:
		
		\[
		12(x+5)=180
		\]
		
		Interpretar \(x\) como número de sectores de riego.
		
	\end{enumerate}
	
	%================================================
	\section*{III. Ecuaciones Fraccionarias}
	
	\begin{enumerate}[label=\textbf{\arabic*.}, start=6]
		
		\item Resolver:
		
		\[
		\frac{x}{4} + 35 = 95
		\]
		
		Interpretar \(x\) como litros de agua distribuidos.
		
		%------------------------------------------------
		
		\item Resolver:
		
		\[
		\frac{2x}{5} - 18 = 42
		\]
		
		Interpretar \(x\) como rendimiento en kg por parcela.
		
		%------------------------------------------------
		
		\item Resolver:
		
		\[
		\frac{x+30}{3}=70
		\]
		
		Interpretar \(x\) como producción de plantas injertadas.
		
	\end{enumerate}
	
	%================================================
	\section*{IV. Ecuaciones Literales}
	
	\begin{enumerate}[label=\textbf{\arabic*.}, start=9]
		
		\item Despejar \(x\):
		
		\[
		P = 15x + 120
		\]
		
		Interpretar \(x\) como superficie cultivada.
		
		%------------------------------------------------
		
		\item Despejar \(r\):
		
		\[
		C = rx + 350
		\]
		
		Interpretar \(r\) como costo por hectárea.
		
		%------------------------------------------------
		
		\item Despejar \(q\):
		
		\[
		T = \frac{q}{4} + 25
		\]
		
		Interpretar \(q\) como cantidad total de fertilizante.
		
	\end{enumerate}
	
	%================================================
	\section*{V. Problemas Contextualizados}
	
	\begin{enumerate}[label=\textbf{\arabic*.}, start=12]
		
		\item Un sistema de riego tecnificado entrega agua a una plantación.
		
		Cada planta requiere 2,8 litros diarios.
		El estanque principal contiene 980 litros,
		pero 140 litros deben reservarse para limpieza.
		
		\begin{enumerate}[label=\alph*)]
			\item Modele la situación con una ecuación fraccionaria.
			\item Determine cuántas plantas pueden regarse.
			\item Analice si el resultado es razonable.
		\end{enumerate}
		
		%------------------------------------------------
		
		\item La producción total de un invernadero se modela mediante:
		
		\[
		P = ax + 500
		\]
		
		donde:
		
		\begin{itemize}
			\item \(P\): producción total
			\item \(a\): rendimiento por hectárea
			\item \(x\): hectáreas cultivadas
		\end{itemize}
		
		\begin{enumerate}[label=\alph*)]
			\item Despeje \(x\).
			\item Si \(P=3200\) y \(a=45\), determine \(x\).
			\item Interprete el resultado.
		\end{enumerate}
		
		%------------------------------------------------
		
		\item Un agricultor compra fertilizante a crédito.
		
		Debe pagar una deuda total de \$1.200.000.
		
		Ya ha abonado \$350.000 y pagará
		el resto en cuotas iguales de \$85.000.
		
		\begin{enumerate}[label=\alph*)]
			\item Plantee la ecuación.
			\item Determine el número de cuotas.
			\item Evalúe si existe consistencia práctica.
		\end{enumerate}
		
	\end{enumerate}
	
	%================================================
	\section*{VI. Patrones e Inconsistencias}
	
	\begin{enumerate}[label=\textbf{\arabic*.}, start=15]
		
		\item Observe las siguientes ecuaciones:
		
		\[
		10x + 100 = 200
		\]
		
		\[
		20x + 100 = 300
		\]
		
		\[
		30x + 100 = 400
		\]
		
		\[
		40x + 100 = 500
		\]
		
		\begin{enumerate}[label=\alph*)]
			\item Resuelva cada ecuación.
			\item Detecte el patrón presente.
			\item Explique la regularidad matemática observada.
		\end{enumerate}
		
		%------------------------------------------------
		
		\item Un estudiante obtiene:
		
		\[
		15x + 60 = 0
		\]
		
		y concluye que:
		
		\[
		x = -4
		\]
		
		siendo \(x\) el número de sacos de fertilizante.
		
		\begin{enumerate}[label=\alph*)]
			\item ¿El resultado algebraico es correcto?
			\item ¿El resultado contextual es válido?
			\item Explique la inconsistencia detectada.
		\end{enumerate}
		
		%------------------------------------------------
		
		\item Analice la siguiente afirmación:
		
		``Si una ecuación tiene solución negativa,
		entonces siempre está mal resuelta.''
		
		\begin{enumerate}[label=\alph*)]
			\item ¿Es verdadera o falsa?
			\item Justifique con un ejemplo agrícola.
		\end{enumerate}
		
	\end{enumerate}
	
	%================================================
	
	\begin{center}
		\Large
		\textbf{Fin de la guía}
	\end{center}
	
\end{document}